\chapter{为什么设备操作需要长期请求对象}

\begin{statebox}
\textbf{项目里程碑 v0.7。} 内核已经重映射 8259A，接收 PIT IRQ0 与键盘 IRQ1，并以
同步 ATA PIO 重读 LBA 0。IRQ 入口和 EOI 解决“事件怎样到达并被确认”，却没有长期
请求身份；寄存器写入也不表示数据已经传输，更不能说明缓冲区何时可以复用。
\end{statebox}

\inputchapterbody{chapters/ch14-device-drivers-dma}
\inputdetail{chapters18/expansions/ch09-queue-reset-deep-dive}
\Needspace{8\baselineskip}
\inputdetail{chapters18/expansions/ch09-request-lifetime-trace}

\begin{problemframe}
\textbf{尚未解决的问题。} 设备路径已经有请求、完成、超时和移除边界，但案例仍只有
Ring 0 主循环，无法让多个用户工作在等待 I/O 时共享处理器。下一章沿 v0.8 与 v0.9
建立用户执行身份、task、runqueue、wait queue 和抢占调度。
\end{problemframe}
